\documentclass[10pt,a4paper]{amsart}
\usepackage[margin=19mm]{geometry}
\usepackage{amsmath,amssymb,mathtools}
\usepackage[scr=rsfs,cal=euler]{mathalfa}
\usepackage{xfrac}
\usepackage[unicode,hidelinks,bookmarks=false]{hyperref}
\newcommand{\ie}{i.e.\ }
\newcommand{\cf}{cf.\ }
\newcommand{\resp}{respectively\ }
\newcommand{\Subsection}{\subsection}
\providecommand{\nobreakdash}{\nobreak\hbox{-}}
\setlength{\emergencystretch}{2em}
\multlinegap0pt
\usepackage{bm}
\usepackage{bbm}
\usepackage{dsfont}
\usepackage{xspace}
\usepackage{cancel}
\usepackage{mathtools}
\usepackage{amsthm}
\makeatletter
\@ifundefined{thm}{\newtheorem{thm}{Theorem}}{}
\makeatother
\newcommand{\ad}{\mathrm{ad}}
\newcommand{\bj}{\mathbf{j}}
\newcommand{\ddu}[2]{\frac{\partial}{\partial u_{#1}^{(#2)}}}
\newcommand{\mbf}[1]{\mathbf{#1}}
\newcommand{\rk}[1]{\mathrm{rank} \left(#1\right)}
\title[On Driftless Systems with m controls and 2m or 2m-1 states t]{On Driftless Systems with m controls and 2m or 2m-1 states that are Flat by Pure Prolongation}
\author{Jean L\'evine, Jaume Franch}
\date{Source: arXiv:2511.06511v1 (CC BY 4.0)}
\begin{document}
\maketitle

\noindent\emph{Source and licence.} On Driftless Systems with m controls and 2m or 2m-1 states that are Flat by Pure Prolongation. Version arXiv:2511.06511v1 (CC BY 4.0), distributed under the Creative Commons Attribution 4.0 International licence, as recorded in the arXiv licence metadata for this exact version and shown on the abstract page https://arxiv.org/abs/2511.06511v1. Full rights notice: licenses/arxiv-2511.06511-CC-BY-4.0.txt.\par\smallskip

\noindent\textit{Editorial scope.} Counted unit: the Thm whose source label is \texttt{3-5theorem} in the article (source0.tex lines 464--479, source label \texttt{3-5theorem}), delivered together with the article's own complete original proof of it (source0.tex lines 481--559, one \texttt{proof} environment of 79 lines). No mathematics is written, repaired or re-derived: statement and proof are the article's own text line for line. Required context retained uncounted: 3inputs (lines 295--297); Deltadef:eq (lines 317--322); CNSP2:thm (lines 329--337); sys5-3:eq (lines 451--453). Every definition, display and label of those lines is retained unchanged. Omitted from this package and not redistributed: 1--294, 298--316, 323--328, 338--450, 454--463, 480--480, 560--1048. Nothing inside a retained excerpt is omitted or altered: every retained body is delivered exactly as the article's lines are written, so no omission receipt of any kind accompanies this package. Citations: the retained excerpts carry no citation command at all, so no bibliography entry of the article is transcribed and no reference is redistributed. Reference adaptation: none was needed and none was made; every reference inside the delivered unit resolves inside it, so no unresolved reference or citation is produced. Printed slips kept literal: no printed word, symbol, hypothesis or number of the counted statement or of its proof was altered. Layout and class: the article's own class is \texttt{article} with packages including graphicx, amssymb, epstopdf, amsmath, amsthm, color; the wrapper is the lane's pinned \texttt{amsart} editorial wrapper at 10pt on A4, which restates the theorem-like environments the retained text needs and adds the article's own macro definitions that the retained text needs (\texttt{\textbackslash ad}, \texttt{\textbackslash bj}, \texttt{\textbackslash ddu}, \texttt{\textbackslash mbf}, \texttt{\textbackslash rk}), with extra wrapper packages (bm, bbm, dsfont, xspace, cancel, mathtools). The delivered numbering is the unit's own. Assets: no retained span contains a figure environment, a graphics-inclusion command, a PDF-page inclusion, a table or a TikZ picture; no image file is copied into this package, the package declares no asset dependency and no image file is redistributed with it. Licence: version arXiv:2511.06511v1 of this article is distributed under the Creative Commons Attribution 4.0 International licence, as recorded in the arXiv licence metadata for this exact version and shown on the abstract page https://arxiv.org/abs/2511.06511v1. A full rights notice is delivered at licenses/arxiv-2511.06511-CC-BY-4.0.txt. Characters: every retained excerpt is pure ASCII, so the delivered engine, XeLaTeX with its default Latin Modern text font, reports no missing character.
\begin{equation}\label{3inputs}
\dot{x} = u_{1}^{(0)}g_{1}(x) + u_{2}^{(0)}g_{2}(x) + u_{3}^{(0)}g_{3}(x)
\end{equation}

\begin{equation}\label{Deltadef:eq}
\begin{aligned}
\Delta_{k}^{(\mbf{j})} &\triangleq \sum_{p=1}^{m} \left\{ \ad_{g_{0}^{(\mbf{j})}}^{l-j_{p}}\ddu{p}{0}  \mid l= j_{p}, \ldots, k\ \right\}\\
\Gamma_{k}^{(\mbf{j})} &\triangleq \bigoplus_{i=1, \ldots, m}\{\ddu{i}{j_{i}-r} \mid r=0, \ldots, k\vee(j_{i}-1))\}
\end{aligned}
\end{equation}

\begin{thm}
\label{CNSP2:thm}
The necessary and sufficient conditions for $P^2$-flatness are:
\begin{itemize}
\item[(i)] $[\Delta_{k}^{(\mbf{j})}, \Delta_{k}^{(\mbf{j})}] \subset \Delta_{k}^{(\mbf{j})}$ and $\rk \Delta_{k}^{(\mbf{j})}$  locally constant for all $k$,
\item[(ii)]$[\Gamma_{k}^{(\mbf{j})}, \Delta_{k}^{(\mbf{j})}] \subset \Delta_{k}^{(\mbf{j})}$ for all $k$,
\item[(iii)] $\exists k_{\star}^{(\mbf{j})} \leq n+\mid \bj \mid$ such that $\rk \Delta_{k}^{(\mbf{j})} = n+m$ for $k\geq k_{\star}^{(\mbf{j})}$.
\end{itemize}
\end{thm}

\begin{equation}\label{sys5-3:eq}
\dot{x} = u_{1}^{(0)}g_{1}(x) + u_{2}^{(0)}g_{2}(x) + u_{3}^{(0)}g_{3}(x)
\end{equation}

\begin{thm}
\label{3-5theorem}
Assume that $\{g_1,g_2,g_3\}$ is not involutive.

\begin{enumerate}
  \item If the largest involutive subdistribution of $\{g_1,g_2,g_3\}$ is $\{g_1,g_2\}$, a necessary and sufficient condition for the system (\ref{sys5-3:eq}) to be flat by pure prolongation is
      \[ H_{1,3}=\left\{g_1,g_2,g_3,[g_{3},g_{1}], [g_{3},g_{2}]\right\} \mbox{ has rank $5$ } \]
      and the minimal prolongation  order is $\mbf{j}= (0,0,1)$. Moreover, the flat outputs are $y_1,y_2,y_3$ such that their differentials annihilate $g_1$ and $g_2$.
  \item If the largest involutive subdistribution of $\{g_1,g_2,g_3\}$ is $\{g_1\}$, a necessary and sufficient condition for the system (\ref{sys5-3:eq}) to be flat by pure prolongation is \\
\centerline{$H'_{1,2} =  \left\{g_{1}, g_{2}, [g_{2},g_{1}], [g_{3},g_{1}]\right\}$ involutive, with $\rk H' _{1,2} = 3$,}\\
\centerline{$\rk H_{2,3} = 5$}\\
\centerline{ with $H_{2,3} = \left\{g_{1}, g_{2}, g_{3}, [g_{2},g_{1}], [g_{3},g_{2}], [g_{3},[g_{3},g_{1}]], [g_{3},[g_{3},g_{2}]] \right\}$,}\\
 and the minimal prolongation order is $\mbf{j}= (0,1,2)$. \\
 Moreover, the flat outputs are $y_1,y_2,y_3$ such that $y_1,y_2$ are independent of the inputs and their differentials annihilate $H'_{1,2}$, while $y_3$ is independent of $u_1$ and $u_2$ and its differential annihilate $g_1$.
\end{enumerate}
\end{thm}

\begin{proof}

\begin{enumerate}
  \item After adding an order one prolongation of $u_3$, the drift and the input vector fields  of the prolonged system are:

\begin{equation}\label{prolvecfields5states:eq}
\begin{aligned}
&g_{0}^{(\mbf{j})}= \left(u_{1}^{(0)}g_{1}(x) + u_{2}^{(0)}g_{2}(x) + u_{3}^{(0)}g_{3}(x)\right) +   u_{3}^{(1)}\ddu{3}{0}\\
&g_{1}^{(\mbf{j})}= \ddu{1}{0}, \qquad g_{2}^{(\mbf{j})}= \ddu{2}{0}, \qquad g_{3}^{(\mbf{j})}= \ddu{3}{1}.
\end{aligned}
\end{equation}

We compute the distributions (\ref{Deltadef:eq}) for $k=0,1,2$:
\begin{equation}
\begin{aligned}
 \Gamma_{0}^{(\mbf{j})} &\triangleq \left\{ \ddu{3}{1} \right\} \\
\Delta_{0}^{(\mbf{j})} &\triangleq \left\{ \ddu{1}{0}  , \ddu{2}{0} \right\}\\
\Gamma_{1}^{(\mbf{j})} &\triangleq \left\{ \ddu{3}{1} \right\} \\
\Delta_{1}^{(\mbf{j})} &\triangleq \left\{ \ddu{1}{0}  , \ddu{2}{0}, \ddu{3}{0}, g_1, g_2 \right\}\\
\Gamma_{2}^{(\mbf{j})} &\triangleq \left\{ \ddu{3}{1} \right\} \\
\Delta_{2}^{(\mbf{j})} &\triangleq \left\{ \ddu{1}{0}  , \ddu{2}{0}, \ddu{3}{0}, g_1, g_2,g_3,[g_0^{(\mbf{j})},g_1], [g_0^{(\mbf{j})},g_2] \right\}
\end{aligned}
\end{equation}

But, due to the involutivity of $H_{0,2}$,
\[ \Delta_{2}^{(\mbf{j})} =\left\{ \ddu{1}{0}  , \ddu{2}{0}, \ddu{3}{0}, g_1, g_2,g_3,[g_3,g_1], [g_3,g_2] \right\} \]
which has rank $8$ thanks to the hypothesis that the rank of $H_{1,3}$ is $5$.

It is a straightforward computation to check the necessary and sufficient conditions for P$^2$-flatness given in theorem \ref{CNSP2:thm}:
\begin{itemize}
\item[(i)] $[\Delta_{k}^{(\mbf{j})}, \Delta_{k}^{(\mbf{j})}] \subset \Delta_{k}^{(\mbf{j})}$ and $\rk \Delta_{k}^{(\mbf{j})}$  locally constant for all $k=0,1,2$,
\item[(ii)]$[\Gamma_{k}^{(\mbf{j})}, \Delta_{k}^{(\mbf{j})}] \subset \Delta_{k}^{(\mbf{j})}$ for all $k=0,1,2$,
\item[(iii)] $\rk \Delta_{2}^{(\mbf{j})} = n+m=8$.
\end{itemize}

Therefore, the system is P$^2$-flat with an order one prolongation of $u_3$ and the flat outputs $y_1,y_2,y_3$ are the solutions of the system of PDE's:

\[ \left< G_{1}^{(\bj)},  dy_{i} \right>= 0 \]

or, equivalently, the flat outputs must be independent of $u_1^{(0)},u_2^{(0)},u_3^{(0)}$ and must be the solution of the system of PDE's

\[ \left< H_{0,2},  dy_{i} \right>= 0 \]

Remark that, by the Frobenius theorem, there exists $y_1,y_2,y_3$ independent solutions of this system of linear partial differential equations since, by assumption, $H_{0,2}$ is an involutive distribution.


\item Let us consider now the case $\mbf{j}=(0,1,2)$. The distributions (\ref{Deltadef:eq}) for $k=0,1,2$ are:
\begin{equation}
\begin{aligned}
\Gamma_{0}^{(\mbf{j})} &= \left\{ \ddu{2}{1},\ddu{3}{2} \right\},\\
\Delta_{0}^{(\mbf{j})} &= \left\{\ddu{1}{0}\right\},\\
\Gamma_{1}^{(\mbf{j})} &= \left\{ \ddu{2}{1},\ddu{3}{2}, \ddu{3}{1} \right\},\\
\Delta_{1}^{(\mbf{j})} &= \left\{\ddu{1}{0}, \ddu{2}{0},  g_{1} \right\},\\
\Gamma_{2}^{(\mbf{j})} &= \Gamma_{1}^{(\mbf{j})}, \\
\Delta_{2}^{(\mbf{j})} &= \left\{\ddu{1}{0}, \ddu{2}{0}, \ddu{3}{0}, g_{1}, g_{2}, [g_{2},g_{1}]\right\}.
\end{aligned}
\end{equation}
are all involutive if, and only if, $H'_{1,2}$ is involutive with $\rk H'_{1,2} = 3$.  Moreover, we have $\rk \Delta_{2}^{(\mbf{j})} = 6$ and we have $\Gamma_{3}^{(\mbf{j})} = \Gamma_{1}^{(\mbf{j})}$ and $\Delta_{3}^{(\mbf{j})}=$
$$\left\{\ddu{1}{0}, \ddu{2}{0}, \ddu{3}{0}, g_{1}, g_{2}, g_{3}, [g_{2},g_{1}], [g_{3},g_{2}], [g_{3},[g_{3},g_{1}]], [g_{3},[g_{3},g_{2}]] \right\}.$$
Thus, $\rk H'_{1,2} = 3$ implies that  $\rk \Delta_{3}^{(\mbf{j})} = 8= n+m$ if $\rk H_{2,3} = 5$. \\
Hence, if $H'_{1,2}$ is involutive, with $\rk H' _{1,2} = 3$, and $\rk H_{2,3} = 5$, the system \eqref{3inputs} is P$^{2}$-flat with minimal prolongation $\mbf{j}=(0,1,2)$. The flat outputs $y_1,y_2$ are the solutions of the system of PDE's:

\[ \left< G_{2}^{(\bj)},  dy_{i} \right>= 0 \]

or, equivalently, the flat outputs $y_1,y_2$ must be independent of $u_1^{(0)},u_2^{(0)},u_3^{(0)}$ and must be the solution of the system of PDE's

\[ \left< H'_{1,2},  dy_{i} \right>= 0 \]

Again, by the Frobenius theorem, there exists $y_1,y_2$ independent solutions of this system of linear partial differential equations since, by assumption, $H'_{1,2}$ is an involutive distribution. Finally, the third flat output $y_3$ can be obtained as an independent function of $y_1,y_2$ satisfying the system of PDE's:

\[ \left< G_{1}^{(\bj)},  dy_{3} \right>= 0 \]

or, equivalently, $y_3$ must be independent of $u_1^{(0)},u_2^{(0)}$ and must be the solution of the system of PDE's

\[ \left< g_1,  dy_{i} \right>= 0 \]

\end{enumerate}

\end{proof}

\end{document}
